\documentclass[10pt,a4paper]{amsart}
\usepackage[margin=19mm]{geometry}
\usepackage{amsmath,amssymb,mathtools}
\usepackage[scr=rsfs,cal=euler]{mathalfa}
\usepackage{xfrac}
\usepackage[unicode,hidelinks,bookmarks=false]{hyperref}
\newcommand{\ie}{i.e.\ }
\newcommand{\cf}{cf.\ }
\newcommand{\resp}{respectively\ }
\newcommand{\Subsection}{\subsection}
\providecommand{\nobreakdash}{\nobreak\hbox{-}}
\setlength{\emergencystretch}{2em}
\multlinegap0pt
\usepackage{amssymb}
\usepackage{mathtools}
\newtheorem{assumption}{Assumption}
\newtheorem{theorem}{Theorem}
\newtheorem{corollary}{Corollary}
\newtheorem{lemma}{Lemma}
\newcommand{\eqnst}[1]{\begin{equation*}#1\end{equation*}}
\title{Scaling and partial universality of the height zero probability in the 2D Abelian sandpile}
\author{Miles W. Elvidge, Antal A. J\'arai}
\date{Source: arXiv:2607.26169v1 (CC BY 4.0)}
\begin{document}
\maketitle

\noindent\emph{Source and licence.} Scaling and partial universality of the height zero probability in the 2D Abelian sandpile. Version arXiv:2607.26169v1 (CC BY 4.0), distributed by arXiv under the Creative Commons Attribution 4.0 International licence, http://creativecommons.org/licenses/by/4.0/, as recorded in the arXiv licence metadata for this exact version and shown on the abstract page https://arxiv.org/abs/2607.26169v1. Full licence notice: \texttt{licenses/arxiv-2607.26169-CC-BY-4.0.txt}.\par\smallskip

\noindent\textit{Editorial scope.} Counted unit: the article's lemma lem:fU>0 (source0.tex lines 1107--1110). The article's complete original proof of it is delivered with it (source0.tex lines 1112--1122, one \texttt{proof} environment of 11 lines). No mathematics is written, repaired or re-derived: the statement and the proof are the article's own lines, in the article's own notation, and no retained line is substituted or re-flowed.  Retained uncounted as the source-local context the proof needs: lines 222--231; lines 988--996; lines 1055--1060; lines 1075--1079.  Nothing was omitted from any retained span, so the package carries no omission record.  No retained line contains a citation, so the wrapper carries no bibliography and no bibliography entry.  No retained line contains a figure environment, an image-inclusion command or drawing code, so no image is delivered and the package declares no asset dependency.  Editorial formatting only, in addition: an AMS \texttt{amsart} preamble that restates the article's own declarations and macros used by the retained lines, namely the environments \texttt{assumption}, \texttt{theorem}, \texttt{corollary}, \texttt{lemma} on separate counters, together with the standard packages those lines need.  The retained environments print in this document as Assumption 1, Theorem 1, Corollary 1, Lemma 1 in order of appearance, whereas the article prints the same items with the numbering of its own full document.  The complete native source of the article as distributed in the arXiv e-print for this exact version is bundled unchanged as \texttt{sources/206244/source0.tex}.
\begin{assumption}
\label{a:U} 
Let $U \subsetneq \mathbb{C}$ be a region (connected open set) such that for each point
$a \in \partial_\infty U$ (considered in the Riemann sphere) there is a continuous simple
curve starting at $a$ and lying outside $U$ (i.e.~there is an injective 
$\gamma: [0,1] \to U^c$ with $\gamma(0) = a$.
It will be convenient for the proof to distinguish the regions we consider as follows. \ \\
(i) $U$ is bounded; \\
(ii) $U$ is unbounded with at least one unbounded boundary component.
\end{assumption}

\begin{theorem}
\label{thm3_6}
    Let \(U, V \subsetneq \mathbb{C}\) be regions satisfying Assumption \ref{a:U}, and assume that
    $\varphi : U \to V$ is a conformal map. Let $z \in U$ and let $w := \varphi(z) \in V$. 
It holds that 
\eqnst{
 f_U(z)
 = |\varphi'(z)|^2 f_V(\varphi(z)). }
\end{theorem}

\begin{corollary}\label{cor3_7}
    Let \(U\subset\mathbb{C}\) be a simply connected region satisfying Assumption \ref{a:U}. Let \(z\in U\).  
Then 
    \[  f_U(z)
        = \frac{2}{(\mathrm{rad}_U(z))^2}.   \]
\end{corollary}

\begin{lemma}
\label{lem:monotone} \ \\
(i) If $z \in U$, we have $f_U(z) \ge 0$. \\
(ii) If $z \in U \subset V$, we have $f_U(z) \ge f_V(z)$.
\end{lemma}

\begin{lemma}
\label{lem:fU>0}
Suppose that $U$ satisfies Assumption \ref{a:U}. If $z \in U$, we have $f_U(z) > 0$.
\end{lemma}

\begin{proof}
First assume that $U$ is bounded.
Take $V \supset U$ to be an open disc in Lemma \ref{lem:monotone}(ii). Then by Corollary \ref{cor3_7}, 
we have $f_V(z) > 0$, and the claim follows from Lemma \ref{lem:monotone}(ii).

In the unbounded case, we can map $U$ conformally to a bounded region satisfying Assumption
\ref{a:U}(i). Indeed, there is a simple curve $\gamma$ in the Riemann sphere $\mathbb{C}_\infty$
starting at $\infty$ such that $U \subset V := \mathbb{C} \setminus \gamma$. Then a conformal
mapping of $V$ to a bounded simply connected region exists, and this maps $U$ conformally 
to a bounded region. Now the claim follows from Theorem \ref{thm3_6}.
\end{proof}

\end{document}
